% Nouveau EGA V (cote 25): archive cote 25, site page 182 (PDF page 183).
% typescript. Draft compared: none.
% AI-transcribed by Claude Opus 5.5 (claude-opus-5-5, adaptive thinking, effort high), Message Batches API, batch ega25b (msgbatch_01A7n5YuXubUPNK1zsepi8RG), 2026-11-13; not yet checked by eye.
\scanpage{25/182}
\mspage{\hand{9}}

terminologie reçue  - je n'ai pas Bourbaki sous la main pour vérifier ce
point- ) devrait s'appeler \uline{déterminant} de la forme quadratique, et non
discriminant. J'aimerais connaitre ton opinion sur cette question; si
tu es d'accord, \unsure{neus|nous} profiterions de l'occasion pour rectifier sur ce
point la terminologie reçue, qui induit en erreur (puisque jusqu'à ces
derniers jours, j'ai sans m'en rendre compte confondu discriminant et
déterminant).

\uline{Remarques} 1\over{\unsure{5}}{6}.12. Lorsqu'on fait opérer le schéma en groupes absolu
$Gl(n)$ \xout{xxxxZ} $= \underline{Aut}(\underline{Z}^n )$ sur $\underline{Q}(\underline{Z}^n)$ , le stabilisateur de la forme
quadratique standart est désigné par $\underline{O}(n)$, et s'appelle le ``\uline{groupe ortho-}%
\uline{gonal absolu}''.( Plus généralement, pour toute forme quadratique $Q$ \hand{$\in$}
\unsure{$\Gamma$|T} $\underline{Sym}^2(\underline{E})$ , on introduit le sous-préschéma en groupes fermé de
$Gl(\underline{E})$ \xout{xxx} stabilisateur de $Q$ , noté $\underline{O}(Q)$ et appelé \uline{schéma en groupes}
\uline{orthogonal relatif à la forme Q}; lorsque $Q$ est isomorphe à l'image
inverse de la forme standart \inserted[typed]{à n variables}\hand{,} $\underline{O}(Q)$ est isomorphe à $O(n)_S$ ). Ceci posé,
il est immédiat que deux points \inserted{$Q,Q'$} du premier membre de (**), à valeurs dans
un $S'$, ont même image dans le deuxième membre, \unsure{sss|ssi} \xout{ilsxsontxcongruxxxie}
\xout{deuxième} on a $u' = uv$ , avec $v \in \underline{O}(n)(S')$; compte tenu de 16.11. (et
avec la terminologie qui sera développée de façon plus détaillée dans
EGA V ou VI\inserted{voir en attendant SGAD \uline{IV}}) on voit donc que les opérations naturelles à droite de
$\underline{O}(n)_S$ sur le premier membre de (**), et la projection (**), font de
$\underline{Isom}(\underline{E}_0,\underline{E})$ un \uline{fibré principal homogène sur} $\underline{Q}(\underline{E})^*$ , \uline{de groupe} $\underline{O}(n)_{\underline{Q}(\underline{E})^*}$ ;
plus particulièrement, \xout{Gl(E$_0$)xxxxGl} $Gl(n)$ est un fibré principal
homogène sur $\underline{Q}(n)^* = \underline{Q}(\underline{Z}^n)^*$, de groupe structural $\underline{O}(n)_{\underline{Q}(n)^*}$ . Il
s'ensuit \xout{xxx} d'abord que (avec les notations de la démonstration de 16.9)
P est en fait un fibré principal homogène sous le groupe $\underline{O}(n)_S$ , associé
\xout{xxxxxiqxxxxx} d'ailleurs canoniquement à la forme $Q$ (de façon fonctorielle
pour les isomorphismes de formes, et de façon compatible avec les change\_%

\margin{\hand{182}}
